\sect{Дискуссия и открытые вопросы}{discussion}
%
В настоящем разделе обсуждаются открытые вопросы и ограничения,
вытекающие из проведённого теоретического анализа, а также
направления возможных исследований, которые выходят за рамки
формальных рассуждений настоящей работы.

\paragraph*{Связь с обзорной таксономией Кима и
коллег.} Обзорная работа Кима с
соавторами~\cite{Kim2026LiveAgents} систематизирует $184$~системы
живых музыкальных агентов в четырёх измерениях~--- Usage Context,
Interaction, Technology, Ecosystem~--- с $31$~детальной осью.
В этой таксономии теоретическое предложение настоящей
работы занимает положение, для которого не нашлось
прямых аналогов: технологически~--- это \textit{гибрид
вероятностной модели и обучаемой политики} (а не
end-to-end нейросеть и не чисто символическая система); по
характеру взаимодействия~--- это \textit{accompaniment с
адаптацией темпа} (а не генерация нового материала и не
free improvisation); по runtime/latency~--- это \textit{soft
real-time с целевым бюджетом $\le 20$\,мс}. Можно
утверждать, что обзор~\cite{Kim2026LiveAgents} косвенно
подтверждает существование пробела в описанной
комбинации факторов и таким образом косвенно
обосновывает потенциальную ценность предлагаемой
архитектуры.

\paragraph*{Перенос методов из астроинформатики: критический
взгляд.} Рассмотренные в~\Secref{ssmm-df} методы SSMM и DF
теоретически могут расширить набор инструментов отслеживания
партитуры. В то же время необходима академическая осторожность
при оценке их вклада. Идея data-driven оценки матрицы переходов
HMM из символьного потока сама по себе не
нова~\cite[\S~III-C]{Rabiner1989}, и преимущества именно
\textit{slotted}-варианта SSMM перед стандартным алгоритмом
Баума-Уэлча в условиях, когда период исполнения известен
заранее (для музыки это так), не очевидны без проверочных
экспериментов. Аналогично DF-признак концептуально близок к
chroma-векторам~\cite{Mueller2007} и \textit{cyclic recurrence
plots}~\cite{MullerFMP2015}, и его методологическое преимущество
сводится к явной двумерной структуре «фаза-амплитуда»;
\textit{количественное} сравнение с уже устоявшимися признаками
является направлением последующей работы. Таким образом, мы
полагаем разумным включить SSMM и DF в перечень кандидатов на
проверку, но не в основной вклад статьи.

\paragraph*{Multiview metric learning и
LM\textsuperscript{3}L.} Естественным
обобщением многомерного признакового пространства, объединяющего
forward-распределение, оценку темпа, SSMM-фичу и DF-фичу
текущего такта, является \textit{multiview metric
learning}~\cite{HuLM3L2014,Bellet2014Metric}. Такая постановка
позволяла бы обучать единую метрику расстояния, по которой
работает $k$-NN-классификатор для задачи перевыравнивания
(re-synchronization) в случае потери трекером уверенности. Тем
не менее формальной публикованной постановки multiview metric
learning \textit{внутри} цикла отслеживания партитуры авторам
не известно; такая разработка~--- предмет потенциальной
самостоятельной работы~\cite{Johnston2020}.

\paragraph*{Обучение по предпочтениям как направление развития.}
В предлагаемой постановке~\Secref{proposal} функция
вознаграждения~\eqref{eq:rl-reward-proposal} определена через
квадратичные и линейные комбинации измеримых временных
расхождений и темповой гладкости. В то же время
\textit{музыкальная} удовлетворительность аккомпанемента не
сводится только к этим характеристикам: интерпретация
динамики, фразировка, артикуляция, баланс между голосами
оркестра~--- всё это эстетические аспекты, плохо описываемые
явной формулой. Принципиальное направление продолжения
работы~--- замена явного вознаграждения на парные сравнения
исполнителей-экспертов (preference learning) по
схеме~\cite{Christiano2017Preference}. Соответствующая Bradley-
Terry форма~\Eqref{rlhf-pref} непосредственно применима к
сегментам исполнения длительностью несколько секунд; сбор
обучающего датасета сравнений у профессиональных
консерваторских музыкантов является самостоятельной
организационной задачей.

\paragraph*{Формулировка через POMDP над belief-состоянием.} Ещё
одно теоретически интересное направление~--- замена локального
forward-распределения на полное \textit{belief-состояние} в
смысле POMDP~\cite{Kaelbling1998POMDP}. В этой постановке
RL-агент действует не над текущей наилучшей оценкой позиции, а
над всем апостериорным распределением, и пространство действий
расширяется до дискретного множества \{«продолжать следить»,
«переключиться на резервный OLTW», «инициировать
ресинхронизацию»\}. Такая постановка теоретически объединяет
управление темпом и управление режимом трекера, но
аналитическое исследование сходимости и стабильности подобной
агентной архитектуры выходит за рамки настоящей работы.

\paragraph*{Этические и художественные аспекты.} В завершение
кратко отметим, что задача автоматического аккомпанемента
имеет художественное измерение, которое нельзя свести к
численным метрикам. Виртуальный оркестр, идеально следующий за
солистом, не заменяет художественного диалога между
живыми музыкантами; уместная задача для систем подобного
рода~--- \textit{обучающая} и \textit{репетиционная}: дать
учащемуся возможность тренироваться с виртуальным
сопровождением при отсутствии живых партнёров, но не
претендовать на полноценное исполнительское взаимодействие. В
обзоре~\cite{Kim2026LiveAgents} эта дискуссия проводится в
контексте измерения «Ecosystem» и заслуживает развёрнутого
обсуждения в отдельной публикации.
